\scanpage{40}
% archive foliation: 39
\deleted{Bien sûr, par définition,}

\uline{Il \unsure{faudr.} expliciter} que pour une catégorie $I$ d'indices, telle que les Inf des parties non vides existent, les foncteurs $X\longrightarrow I$ sont en corr. $1$-$1$ avec les filtrations \deleted{\unsure{croiss.}} applications croissantes
\[
\begin{aligned}I&\longrightarrow\operatorname{SsCatOuv}(X)\\i&\longmapsto X_i\end{aligned}
\]
qui commutent aux dits Inf, et telles que l'on ait $X=\bigcup X_i$.
\margin{cette dernière condition équivaut aussi à $\varinjlim X_i\simeq X$.}

Ceci posé, on a
\[
X_{/i}=X_i
\]
\deleted{d'où \unsure{l'explicitation}, pour $f$}
\[
f_i:X_i\mathbin{-}Y_i
\]
s'identifie à
\[
f_{/i}:X_{/i}\mathbin{-}Y_{/i}
\]
donc on gagne par L3.

\margin{$\vert$}
\uline{Corollaire \unsure{1|!}} Idem dans un $\widehat A$, en utilisant le fait que $i_A$ commute aux \uline{$\varinjlim$} quelconques et \deleted{aux} produits fibrés.

\margin{\uline{Cor?} Si les $X_i$ sont \inserted{$\underline W$-}asphériques et $I$ connexe, alors $X$ est $\underline W$-asphérique (On prend $Y_i=\unsure{e_{(I_i)}|e_i}$).}

\uuline{Thm} L'ens. $\underline W$ est saturé

\uline{Cor} Pour $\forall A$, $\underline W_A$ est saturé.

\deleted{
\uline{Dém} Soit $f:X\mathbin{-}Y$ tel que $\unsure{\sigma|\gamma}(f)$ iso, d'où
\diagram{\begin{tikzcd}[column sep=small,row sep=small]
X\arrow[dr,"f"]&&X\arrow[d,"\alpha\in\underline W"']\arrow[drr,"f"]\\
&Y\arrow[r,"g"]\arrow[l,phantom,"\unsure{\exists}"]&X'\arrow[drr,"f'"']&&Y\arrow[d,"\beta\in\underline W"]\\
&&&&Y'
\end{tikzcd}}
\deleted{avec $\unsure{\sigma|\gamma}(gf)=\unsure{\sigma|\gamma}(\alpha)$, et}
avec $\unsure{\sigma|\gamma}(gf)=\unsure{\sigma|\gamma}(\alpha)$, $\unsure{\sigma|\gamma}(f'g)=\unsure{\sigma|\gamma}(\beta)$, $\beta f\simeq f'\alpha$.

Quitte à changer $\alpha$ en composé avec $X'\mathbin{-}X'_1$, et item pour $Y'$ avec $Y'\mathbin{-}Y'_1$, OPS $gf\sim\alpha$, $f'g\sim\beta$, \inserted{donc $gf\in\underline W$, $f'g\in\underline W$} \deleted{et pr (quitte à changer encore $\alpha$ et $\beta$ par $gf$ et $f'g$) \unsure{qui sont égales}}
\diagram{\begin{tikzcd}
X\arrow[r,"f"]\arrow[d,"\alpha"']&Y\arrow[d,"\beta"]\arrow[dl,"g"']\\
X'\arrow[r,"f'"']&Y'
\end{tikzcd}}
\[
\underbrace{\begin{gathered}\alpha,\beta\in W\\gf=\alpha,\quad f'g=\beta\\(\text{quitte à \unsure{factoriser}})\end{gathered}}
\]
\margin{\diagram{\begin{tikzcd}[column sep=small]
X\arrow[r,"f"']&Y\arrow[r,"g"']&Y'\arrow[r,"f'"']&X'
\end{tikzcd}} $gf\in\underline W$, $f'g\in\underline W$}

On pourrait d'ailleurs supposer que $f$ soit imm. ouverte, puis (quitte à factoriser \deleted{$g$, puis $f'$}) que $f,g,f'$ \unsure{sont|soient} des imm. ouvertes.
\unsure{\diagram{\begin{tikzcd}[row sep=large,column sep=large]
X\arrow[rr,hook]&&Y\arrow[dl,"g_1"']\\
&X'_1\arrow[dl,shift right]\arrow[dr,hook]\\
X'\arrow[ur,shift right]\arrow[rr]&&Y'
\end{tikzcd}}}
}
